\weeknum{6}
\topic[Magnetics]{Magnetics: Anomalies and Susceptibility}

\workshopstart

\vspace{-0.75em}
\noindent\textbf{Duration:} 3 hours \hfill \textbf{Setting:} Workshop room (disc magnets, iron filings, susceptibility meter)

\section*{Preparation}

\begin{description}
  \item[Pre-reading:] Course Notes, Ch. 8
  \item[Lecture videos:]
  \item[Notes:]
\end{description}

\section*{Purpose}

This workshop builds physical intuition for magnetic fields and their geological sources. Using dipole demonstrations, rock susceptibility measurements, and a case study from Antarctica, students contrast magnetics with gravity (vector nature, polarity, temperature dependence) and connect magnetic structure to crustal composition and thermal state.

\section*{Learning Outcomes}

By the end of this workshop, students should be able to:
\begin{itemize}
  \item Visualise magnetic fields generated by dipoles and combinations of sources
  \item Relate magnetic susceptibility to rock type and geological process
  \item Predict magnetic anomaly shapes from simple geological geometries
  \item Contrast magnetic anomalies with gravity anomalies for the same structures
  \item Explain the concept of Curie depth and its relationship to thermal structure
  \item Link source depth to anomaly wavelength and spectral content
\end{itemize}

\section*{Session Flow (180 min)}

\timeslot{Introduction and framing}{10}
\begin{itemize}
  \item How magnetics differs from gravity (vector nature, polarity, temperature dependence)
  \item Overview of the three activities and how they connect
\end{itemize}

\timelineitem{Magnetic fields from dipoles}{25}
\begin{itemize}
  \item Iron filings and disc magnets
  \item Vertical vs horizontal dipole orientation
  \item Stacked magnets and multiple sources
\end{itemize}

\noindent\textbf{Learning emphasis:}
\begin{itemize}
  \item Dipole field geometry
  \item Orientation-dependent symmetry and polarity
  \item Intuition for induced magnetisation
\end{itemize}

\timelineitem{Whole-class demo: Magnetic susceptibility of rocks}{20}
\begin{itemize}
  \item Susceptibility measurements of igneous, metamorphic, and sedimentary rocks
  \item Effects of alteration, metamorphism, and weathering
\end{itemize}

\noindent\textbf{Learning emphasis:}
\begin{itemize}
  \item Susceptibility as a material property
  \item Why similar compositions can have very different magnetic behaviour
\end{itemize}

\timelineitem{Magnetic anomalies vs gravity}{40}
\begin{itemize}
  \item Reasoning-based sketching of anomalies for idealised geometries
  \item Direct comparison with gravity anomalies
\end{itemize}

\timeslot{Break}{10}

\timelineitem{Curie depth and magnetic crust}{55}
\begin{itemize}
  \item Curie temperature and magnetic crustal thickness
  \item Depth--wavelength relationships
  \item Antarctica case study
  \item Back-of-the-envelope heat-flow estimates
\end{itemize}

\timeslot{Synthesis and wrap-up}{20}
\begin{itemize}
  \item What magnetics can and cannot tell us
  \item Connection to upcoming Fourier/spectral methods
\end{itemize}

\notepage
\notepage

\subimport{activities/}{activity_magnetic_fields}
\subimport{activities/}{activity_rock_magnetism}
\subimport{activities/}{activity_magnetic_structures}
\subimport{activities/}{activity_curie_depth}